기룡이는 힘들게 손에 넣은 두쫀쿠를 친구들과 나눠 먹기 위해 학교로 갔다. 
하지만 기룡이가 두쫀쿠를 가졌다는 소문이 퍼지며 두쫀쿠를 노리고 사람들이 기룡이에게 달려들었다!!!

기룡이는 어쩔 수 없이 자신을 포함하여 $N$명의 사람들이 모두 두쫀쿠를 맛볼 수 있도록 두쫀쿠를 꾸욱 눌러 피자 같은 원판 모양으로 만든 후, 
두쫀쿠를 $N$개의 조각으로 잘라 나눠 먹으려 한다. 
한 번의 칼질은 두쫀쿠 원판의 둘레 위 서로 다른 두 점을 잇는 하나의 선분으로 이루어진다. 칼질의 위치는 자유롭게 정할 수 있다.

여러 번 칼질하기 귀찮았던 기룡이는 두쫀쿠를 $N$개의 조각으로 나누기 위한 칼질의 횟수를 최소화하고자 한다. 각 조각의 크기가 같을 필요는 없고, 단순히 조각만 $N$개이면 된다.

$N$개의 조각을 만들기 위해선 최소 몇 번의 칼질을 해야 할지 구해보자.